\documentclass[10pt,a4paper]{amsart}
\usepackage[margin=19mm]{geometry}
\usepackage{amsmath,amssymb,mathtools}
\usepackage[scr=rsfs,cal=euler]{mathalfa}
\usepackage{xfrac}
\usepackage[unicode,hidelinks,bookmarks=false]{hyperref}
\newcommand{\ie}{i.e.\ }
\newcommand{\cf}{cf.\ }
\newcommand{\resp}{respectively\ }
\newcommand{\Subsection}{\subsection}
\providecommand{\nobreakdash}{\nobreak\hbox{-}}
\setlength{\emergencystretch}{2em}
\multlinegap0pt
\usepackage{mdframed}
\usepackage{mdframed}
\usepackage{mdframed}
\usepackage{mdframed}
\usepackage{tikz}
\usepackage{amssymb}
\usepackage{float}
\usepackage{xcolor}
\usepackage{tcolorbox}
\usepackage{mdframed}
\usepackage{algorithm}\usepackage{algpseudocode}
\newtheorem{definition}{Definition}
\newtheorem{assumption}{Assumption}
\newtheorem{theorem}{Theorem}
\newcommand{\mubar}{\bar{\mu}}
\title{Finite Sample Smeariness on Spheres and Modulation-aware Tests}
\author{Susovan Pal}
\date{Source: arXiv:2603.20974v3 (CC BY 4.0)}
\begin{document}
\maketitle

\noindent\emph{Source and licence.} Finite Sample Smeariness on Spheres and Modulation-aware Tests. Version arXiv:2603.20974v3 (CC BY 4.0), distributed by arXiv under the Creative Commons Attribution 4.0 International licence, http://creativecommons.org/licenses/by/4.0/, as recorded in the arXiv licence metadata for this exact version and shown on the abstract page https://arxiv.org/abs/2603.20974v3. Full licence notice: \texttt{licenses/arxiv-2603.20974-CC-BY-4.0.txt}.\par\smallskip

\noindent\textit{Editorial scope.} Counted unit: the article's theorem thm:dim\_monotone\_modulation\_rot (source0.tex lines 895--924). The article's complete original proof of it is delivered with it (source0.tex lines 926--1019, one \texttt{proof} environment of 94 lines). No mathematics is written, repaired or re-derived: the statement and the proof are the article's own lines, in the article's own notation, and no retained line is substituted or re-flowed.  Retained uncounted as the source-local context the proof needs: lines 415--418; lines 606--618; lines 635--643; lines 713--719; lines 841--862.  Nothing was omitted from any retained span, so the package carries no omission record.  No retained line contains a citation, so the wrapper carries no bibliography and no bibliography entry.  No retained line contains a figure environment, an image-inclusion command or drawing code, so no image is delivered and the package declares no asset dependency.  Editorial formatting only, in addition: an AMS \texttt{amsart} preamble that restates the article's own declarations and macros used by the retained lines, namely the environments \texttt{definition}, \texttt{assumption}, \texttt{theorem} on separate counters, together with the standard packages those lines need.  The retained environments print in this document as Definition 1, Assumption 1, Theorem 1 in order of appearance, whereas the article prints the same items with the numbering of its own full document.  The complete native source of the article as distributed in the arXiv e-print for this exact version is bundled unchanged as \texttt{sources/196902/source0.tex}.
\begin{definition}[\textbf{Asymptotic modulation}]
\label{def:asymptotic_modulation}
Let \(\mubar\) be a unique Fr\'echet mean of \(X\) on \(M\), and let \(\hat\mu_n\) be sample Fr\'echet means. Set \(V:=\mathbb E[d^2(\mubar,X)]\), \(V_n:=\mathbb E[d^2(\mubar,\hat\mu_n)]\), and \(m_n:=nV_n/V\) for \(V>0\).  If \(X\) is non-smeary and \(\lim_{n\to\infty}m_n\) exists, we call \(m_\infty:=\lim_{n\to\infty}m_n\) the \emph{asymptotic modulation}.
\end{definition}

\begin{equation}
\label{eqn:Hessian-general-density}
\boxed{
    H
    =
    2\int_0^\pi\!\!\int_{\mathbb S^{m-1}}
    \left[
    \frac{R}{\tan R}I_m+
    \left(1-\frac{R}{\tan R}\right)\Theta\Theta^{\mathsf T}
    \right]
    \rho(R,\Theta)(\sin R)^{m-1}\,d\Theta\,dR .
}
\end{equation}

\begin{equation}
\label{eqn:Hess}
\boxed{
    H
    =
    \frac{2\operatorname{vol}(\mathbb S^{m-1})}{m}
    \left(\int_0^\pi \rho(R)(\sin R)^{m-1}b_m(R)\,dR\right)I_m .
}
\end{equation}

\begin{assumption}[\textbf{Uniform moment bound}]
\label{assump:momnt-bound}
Let \(\widehat\mu_n\) be a measurable sequence of intrinsic sample Fr\'echet means and set \(U_n:=\exp_N^{-1}(\widehat\mu_n)\).  We assume that there exists \(\delta>0\) such that
\[
    \sup_{n\ge1}\mathbb E\bigl[\|\sqrt n\,U_n\|^{2+\delta}\bigr]<\infty .
\]
\end{assumption}

\begin{definition}[\textbf{Dimensionally comparable rotationally symmetric family}]
\label{def:dimensionally-comparable-distributions}
Let \(N\) denote the north pole.  Fix \(\varepsilon\in[0,\frac{\pi}{2})\) and set \(R_\varepsilon:=\frac{\pi}{2}-\varepsilon\).  For each \(m\ge2\), let \(\mu^{(m)}\) be a probability measure on \(\mathbb S^m\), rotationally symmetric about \(N\), with
\[
\operatorname{supp}\mu^{(m)}\subset \overline B(N;R_\varepsilon).
\]
We say that \(\{\mu^{(m)}\}_{m\ge2}\) is a \emph{dimensionally comparable rotationally symmetric family} if there exists a measurable function \(w:[0,R_\varepsilon]\to[0,\infty)\) such that, for every \(m\ge2\),
\[
0<Z_m:=\int_{0}^{R_\varepsilon} w(R)(\sin R)^{m-1}\,dR<\infty,
\]
and, in normal polar coordinates \(v=R\Theta\in T_N\mathbb S^m\), the density of \(\mu^{(m)}\) with respect to polar volume measure has the form
\[
\rho^{(m)}(R,\Theta)=\rho_m(R)\,\varphi_m(\Theta),
\qquad
\rho_m(R):=\frac{w(R)}{Z_m},
\]
where \(\varphi_m\) is the uniform probability density on \(\mathbb S^{m-1}\).  Equivalently, the induced radial probability measure on \([0,R_\varepsilon]\) is
\[
\nu_m(dR)=\frac{w(R)(\sin R)^{m-1}}{\int_{0}^{R_\varepsilon} w(U)(\sin U)^{m-1}\,dU}\,dR.
\]
Thus the family is generated by a single dimension-independent unnormalized radial weight \(w\), while the dependence on \(m\) enters only through the Jacobian factor \((\sin R)^{m-1}\) and the normalizing constant \(Z_m\).
\end{definition}

\begin{theorem}[\textbf{Dimension monotonicity of asymptotic modulation for dimensionally comparable rotationally symmetric families}]
\label{thm:dim_monotone_modulation_rot}
Fix \(\varepsilon\in[0,\frac{\pi}{2})\) and set \(R_\varepsilon:=\frac{\pi}{2}-\varepsilon\).  For each integer \(m\ge2\), let \(\mu^{(m)}\) be an absolutely continuous probability measure on \(\mathbb S^m\) with unique intrinsic Fr\'echet mean \(N\), and assume that \(\{\mu^{(m)}\}_{m\ge2}\) is a dimensionally comparable rotationally symmetric family about \(N\), supported in \(\overline B(N;R_\varepsilon)\), in the sense of Definition~\ref{def:dimensionally-comparable-distributions}.  Assume further that the moment condition \ref{assump:momnt-bound} holds for each \(m\).  Let \(X^{(m)}\sim\mu^{(m)}\), set \(\Sigma_m:=\operatorname{Cov}[\exp_N^{-1}(X^{(m)})]\), and let \(H_m:=D^2\widetilde F_m(0)\).  Define the dimension-indexed asymptotic modulations, cf. Definition~\ref{def:asymptotic_modulation}, by
\[
m_\infty^{(m)}
:=
\frac{\operatorname{trace}\!\bigl(4H_m^{-1}\Sigma_m H_m^{-1}\bigr)}{\operatorname{trace}(\Sigma_m)}.
\]
Then:
\begin{enumerate}
\item The sequence \(\bigl(m_\infty^{(m)}\bigr)_{m\ge2}\) is strictly increasing with dimension:
\[
m_\infty^{(m+1)} > m_\infty^{(m)}
\qquad (m\ge2).
\]

\item If the common radial weight \(w\) from Definition~\ref{def:dimensionally-comparable-distributions} satisfies, for every \(\delta>0\),
\[
\int_{R_\varepsilon-\delta}^{R_\varepsilon} w(R)\,dR > 0,
\]
then, in the extended real line,
\[
\lim_{m\to\infty} m_\infty^{(m)}
=
\frac{1}{\bigl(R_\varepsilon \cot R_\varepsilon\bigr)^2}
=
\frac{1}{\Bigl(\left(\frac{\pi}{2}-\varepsilon\right)\cot\!\left(\frac{\pi}{2}-\varepsilon\right)\Bigr)^2}.
\]
\end{enumerate}
\end{theorem}

\begin{proof}
\textbf{Proof of (1).}
By Definition~\ref{def:dimensionally-comparable-distributions}, there exists a measurable function \(w:[0,R_\varepsilon]\to[0,\infty)\) such that \(Z_m:=\int_0^{R_\varepsilon}w(R)(\sin R)^{m-1}\,dR\in(0,\infty)\), \(\rho_m(R)=w(R)/Z_m\), and \(\rho^{(m)}(R,\Theta)=\rho_m(R)\varphi_m(\Theta)\), where \(\varphi_m\) is the uniform probability density on \(\mathbb S^{m-1}\).  Rotational symmetry gives
\[
\int_{\mathbb S^{m-1}}\Theta\Theta^\top\,\varphi_m(\Theta)\,d\Theta=\frac{1}{m}I_m.
\]
Using the Hessian formula \eqref{eqn:Hessian-general-density} with this normalized angular density, equivalently the rotationally symmetric reduction \eqref{eqn:Hess} with the angular volume factor cancelled by \(\varphi_m\), we obtain
\[
H_m
=
\frac{2}{m}\left(\int_0^{R_\varepsilon}\rho_m(R)(\sin R)^{m-1}b_m(R)\,dR\right)I_m.
\]
Since \(b_m(R)=1+(m-1)R\cot R\), this becomes \(H_m=\lambda_m I_m\), where
\[
s_m:=
\int_0^{R_\varepsilon}(R\cot R)\,\rho_m(R)(\sin R)^{m-1}\,dR
=
\frac{\int_0^{R_\varepsilon}(R\cot R)\,w(R)(\sin R)^{m-1}\,dR}{Z_m},
\qquad
\lambda_m:=2\left(s_m+\frac{1-s_m}{m}\right).
\]
Since \(0\le R\cot R<1\) on \((0,\pi/2]\) and \(Z_m\in(0,\infty)\), we have \(0\le s_m<1\).  Hence \(\lambda_m>0\), so \(H_m\) is positive definite.  The trace formula in Definition~\ref{def:asymptotic_modulation} therefore gives
\[
m_\infty^{(m)}
=
\frac{\operatorname{trace}\!\bigl((4/\lambda_m^2)\Sigma_m\bigr)}{\operatorname{trace}(\Sigma_m)}
=
\frac{4}{\lambda_m^2}.
\]
Thus \(m_\infty^{(m+1)} > m_\infty^{(m)}\) is equivalent to \(\lambda_{m+1}<\lambda_m\).

We next prove that \(s_{m+1}<s_m\).  Define the probability measure \(\nu_m\) on \([0,R_\varepsilon]\) by
\[
\nu_m(dR):=\frac{w(R)(\sin R)^{m-1}}{Z_m}\,dR.
\]
Let \(c_m:=\int_0^{R_\varepsilon}\sin R\,\nu_m(dR)=Z_{m+1}/Z_m>0\).  Then
\[
\frac{\sin R}{c_m}\,\nu_m(dR)=\nu_{m+1}(dR).
\]
Therefore
\[
s_{m+1}
=
\frac{1}{c_m}\int_0^{R_\varepsilon}(R\cot R)\sin R\,\nu_m(dR).
\]
Subtracting \(s_m=\int_0^{R_\varepsilon}(R\cot R)\nu_m(dR)\), and using the standard covariance symmetrization identity, gives
\[
s_{m+1}-s_m
=
\frac{1}{2c_m}\int_0^{R_\varepsilon}\int_0^{R_\varepsilon}
\bigl(R\cot R-U\cot U\bigr)(\sin R-\sin U)\,\nu_m(dR)\,\nu_m(dU).
\]
The map \(R\mapsto R\cot R\) is strictly decreasing on \((0,\pi/2]\), while \(R\mapsto\sin R\) is strictly increasing on \([0,\pi/2]\).  Hence the integrand is everywhere nonpositive.  Since \(\nu_m\) is absolutely continuous with respect to Lebesgue measure and is not supported on a single point, the inequality is strict.  Thus \(s_{m+1}<s_m\).

Finally,
\[
\lambda_{m+1}<\lambda_m
\;\Longleftrightarrow\;
m^2(s_{m+1}-s_m)<1-s_m.
\]
Since \(s_{m+1}-s_m<0\) and \(s_m<1\), this inequality holds.  Therefore \(\lambda_{m+1}<\lambda_m\), and hence \(m_\infty^{(m+1)}>m_\infty^{(m)}\).

\textbf{Proof of (2).}
Assume that for every \(\delta>0\), \(\int_{R_\varepsilon-\delta}^{R_\varepsilon}w(R)\,dR>0\).  Fix \(\delta\in(0,R_\varepsilon)\) and write \(q_\delta:=\sin(R_\varepsilon-\delta)\) and \(\beta_\delta:=\int_{R_\varepsilon-\delta/2}^{R_\varepsilon}w(R)\,dR>0\).  Then
\[
\int_0^{R_\varepsilon-\delta}w(R)(\sin R)^{m-1}\,dR
\le
q_\delta^{\,m-2}\int_0^{R_\varepsilon-\delta}w(R)\sin R\,dR,
\]
while
\[
Z_m\ge \beta_\delta\,\sin(R_\varepsilon-\delta/2)^{m-1}.
\]
Since \(q_\delta<\sin(R_\varepsilon-\delta/2)\), it follows that
\[
\nu_m([0,R_\varepsilon-\delta])\to0,
\qquad
\nu_m([R_\varepsilon-\delta,R_\varepsilon])\to1.
\]
As \(R\mapsto R\cot R\) is continuous on \([0,R_\varepsilon]\), we obtain
\[
s_m=\int_0^{R_\varepsilon}(R\cot R)\,\nu_m(dR)
\longrightarrow
R_\varepsilon\cot R_\varepsilon.
\]
Therefore \(\lambda_m=2(s_m+(1-s_m)/m)\to 2R_\varepsilon\cot R_\varepsilon\), and consequently
\[
\lim_{m\to\infty}m_\infty^{(m)}
=
\lim_{m\to\infty}\frac{4}{\lambda_m^2}
=
\frac{1}{\bigl(R_\varepsilon\cot R_\varepsilon\bigr)^2}.
\]
\end{proof}

\end{document}
